%==============================================================================
% ARC 500 - Programming with Python and Generative AI
% Syllabus, Fall 2026
% Junho Chun, PhD - School of Architecture, Syracuse University
%
% Structure and university-policy sections follow the ARC 612 syllabus format.
% Course-specific content is reconciled against the course's own single source
% of truth: ARC500_2026_Semester_Restructure_Weeks01-15.txt (master schedule).
%
% Compile with: pdflatex ARC500_Syllabus_FS26.tex   (run twice for the TOC-free
% cross-references and page totals to settle)
%
% ---------------------------------------------------------------------------
% Official student syllabus, reconciled August 20, 2026 against the course
% master schedule, the Fall 2026 University calendar, current University
% syllabus requirements, and the ARC 612 Fall 2026 common-policy source.
% ---------------------------------------------------------------------------
%==============================================================================
\documentclass[11pt]{article}

\usepackage[utf8]{inputenc}
\usepackage[T1]{fontenc}
\usepackage{lmodern}
\usepackage[margin=1in,headheight=14pt]{geometry}
\usepackage{microtype}
\usepackage{xcolor}
\usepackage{fancyhdr}
\usepackage{booktabs}
\usepackage{longtable}
\usepackage{array}
\usepackage{tabularx}
\usepackage{enumitem}
\usepackage{titlesec}
\usepackage{ragged2e}
\usepackage[hidelinks]{hyperref}

%------------------------------------------------------------------ colors
\definecolor{suorange}{HTML}{D44500}
\definecolor{rulegray}{HTML}{9AA0A6}

%------------------------------------------------------------- course facts
\newcommand{\CourseCredit}{3 credit hours}
\newcommand{\OfficeHours}{Monday and Wednesday, 10:30 am--12:00 pm, and by appointment}
\newcommand{\DropDeadline}{Monday, September 14, 2026}
\newcommand{\WithdrawalDeadline}{Friday, November 20, 2026}
\newcommand{\ProgramRole}{ARC 500 is a 500-level Selected Topics course offered for
  graduate-level credit. This section is a three-credit elective. Its application to a
  student's B.Arch., M.Arch., or M.S. degree requirements depends on the student's program
  and any required faculty or advisor approval. Students should confirm degree applicability
  with their academic advisor.}

%------------------------------------------------------------------ headings
\titleformat{\section}{\normalfont\Large\bfseries\color{suorange}}{}{0pt}{}
\titlespacing*{\section}{0pt}{14pt}{6pt}
\titleformat{\subsection}{\normalfont\large\bfseries}{}{0pt}{}
\titlespacing*{\subsection}{0pt}{11pt}{4pt}

%------------------------------------------------------------------ header/footer
\pagestyle{fancy}
\fancyhf{}
\fancyhead[L]{\small\itshape School of Architecture \textbar{} Syracuse University \textbar{} ARC 500}
\fancyfoot[C]{\small\thepage}
\renewcommand{\headrulewidth}{0.4pt}
\renewcommand{\footrulewidth}{0pt}

\setlist[itemize]{leftmargin=1.4em,itemsep=2pt,topsep=3pt,parsep=0pt}
\setlist[enumerate]{leftmargin=1.6em,itemsep=2pt,topsep=3pt,parsep=0pt}

\newcommand{\msq}{\,m\textsuperscript{2}}
\newcommand{\eui}{kWh/m\textsuperscript{2}/yr}

\begin{document}

%==============================================================================
\thispagestyle{plain}
\begin{center}
  {\LARGE\bfseries ARC 500}\\[5pt]
  {\LARGE\bfseries Programming with Python and Generative AI}\\[5pt]
  {\large Scripting, Data Analysis, Visualization, Problem-Solving, and Optimization}\\[8pt]
  {\Large\bfseries\color{suorange} Syllabus \textbar{} Fall 2026}\\[6pt]
  {\normalsize School of Architecture, Syracuse University}
\end{center}
\vspace{4pt}
{\color{rulegray}\hrule height 1pt}
\vspace{10pt}

%==============================================================================
\section{Contact Information}

\begin{tabularx}{\textwidth}{@{}l X@{}}
  \textbf{Instructor}      & Junho Chun, PhD \\
  \textbf{Office Location} & Slocum Hall 306A \\
  \textbf{Email}           & \href{mailto:jchun04@syr.edu}{jchun04@syr.edu} \\
  \textbf{Phone}           & 315.443.2808 \\
  \textbf{Office Hours}    & \OfficeHours{}. If my office door is open, please feel free to
                             drop by anytime. \\
\end{tabularx}

%==============================================================================
\section{Course Information}

\begin{tabularx}{\textwidth}{@{}l X@{}}
  \textbf{Lecture / Studio} & Monday and Wednesday, 12:45 pm to 2:05 pm \\
  \textbf{Location}         & Slocum Hall 307 \\
  \textbf{Prerequisite}     & None. This course is designed for beginners with no prior
                              coding experience. \\
  \textbf{Credit}           & \CourseCredit{} \\
  \textbf{Required}         & A laptop, brought to \emph{every} class meeting (see
                              \textit{Required Laptop and Software} below) \\
\end{tabularx}

%==============================================================================
\section{Course Objectives}

This course introduces Python programming from the ground up, integrating generative AI as a
practical tool for learning and developing code. Students begin with core programming
concepts---syntax, data types, logic, and functions---and progressively apply these skills to
filter and screen data, create visualizations, solve optimization problems, and explore
introductory machine learning. The emphasis is on building genuine programming literacy:
students learn to write, read, debug, and extend code, with AI serving as an accelerator
rather than a replacement.

Beyond writing scripts, the course provides a foundation in the theoretical principles that
underpin computational problem-solving. Students explore the conceptual and mathematical
foundations of key topics, including optimization theory (objective functions and
constraints), machine learning fundamentals (classification, regression, and model
evaluation), and data analysis methodology. This dual emphasis ensures that students not only
know \emph{how} to implement solutions in code but also understand \emph{why} those methods
work and \emph{when} to apply them. Through a hands-on approach combining instruction and
practice, the course prioritizes professional application over pure theoretical computer
science, ensuring students gain experience that is immediately applicable across disciplines.

As generative AI reshapes the landscape of programming, the primary objective of this course
shifts away from the memorization of syntax. Rather, the curriculum is designed to cultivate
computational thinking, problem formulation, and the evaluation of AI-generated outcomes.
Students who integrate core programming fundamentals with AI literacy will be equipped with
the resilience and adaptability necessary to navigate and leverage future technological
advancements.

%==============================================================================
\section{Learning Goals}

After taking this course, students will be able to:

\begin{itemize}
  \item Write, run, read, trace, test, and debug Python scripts that use core values and types,
        expressions, conditionals, loops, lists, dictionaries, and functions; classify errors
        as syntax, runtime, or logic errors and explain every retained line of code.
  \item Build reproducible NumPy, pandas, and GeoPandas workflows that document provenance and
        units; inspect data quality; clean, group, aggregate, and export data; and justify
        missing-data and outlier decisions.
  \item Select an appropriate visual encoding and produce labeled, accessible, and
        interpretable static, animated, and spatial evidence using Matplotlib and GeoPandas.
  \item Formulate design decisions as bounded optimization problems; implement linear,
        nonlinear gradient-based, and heuristic methods; compare methods; analyze sensitivity;
        and independently verify feasibility and constraints.
  \item Fit and evaluate supervised prediction models against baselines, diagnose limitations
        and leakage, integrate an evaluated model with optimization in a
        predict--optimize--decide workflow, and disclose and verify all generative-AI
        assistance used in the process.
\end{itemize}

%==============================================================================
\section{Course Description}

An introductory, application-focused course in Python and responsible generative-AI-assisted
coding. Students build reproducible workflows for data analysis, visualization, spatial
evidence, optimization, and evaluated machine learning, culminating in a defensible
predict--optimize--decide design recommendation.

%==============================================================================
\section{Course Delivery Format}

The course is delivered in an in-person format, and the lectures will not be recorded. Each
week has two meetings with distinct roles:

\begin{itemize}
  \item \textbf{Meeting A (Monday)} introduces the week's concepts and reasoning, opening with
        a 3--5 minute debrief of the most common error from the previous week's submission.
  \item \textbf{Meeting B (Wednesday)} is a hands-on Spyder studio. Students build that week's
        deliverable in class, with the instructor in the room. \textbf{This is why a laptop is
        required at every meeting.}
\end{itemize}

Three weeks depart from that two-meeting pattern, because of the academic calendar:

\begin{itemize}
  \item \textbf{Week 3} --- Monday, September 7 is Labor Day. Meeting A is replaced by a short
        asynchronous primer (approximately 40 minutes) to be completed before that Wednesday.
  \item \textbf{Week 8} --- Monday, October 12 is Fall Break. There is no Meeting A and no
        substitute is assigned; that week's only work is finishing Project 1.
  \item \textbf{Week 15} --- Monday, December 7 is the only meeting. December 9 is a reading
        day, so the capstone clinic, synthesis, and in-class presentations all occur in that
        single session.
\end{itemize}

\newpage
The University's Thanksgiving Break runs Sunday, November 22 through Sunday, November 29 and
removes both ARC 500 meetings that week. It is treated as an unnumbered break between Weeks 13
and 14, so the numbered sequence stays 1--15.

\subsection*{Key Fall 2026 Dates}

\begin{tabularx}{\textwidth}{@{}l X@{}}
  \textbf{August 24}   & First day of classes and first ARC 500 meeting \\
  \textbf{August 31}   & Add deadline \\
  \textbf{September 14}& Academic/financial drop, grading-option, and religious-observance
                          notification deadline \\
  \textbf{October 12--13} & Fall Break; ARC 500 does not meet Monday, October 12 \\
  \textbf{November 20} & University withdrawal deadline \\
  \textbf{November 22--29} & Thanksgiving Break; no class November 23 or 25 \\
  \textbf{December 7}  & Final ARC 500 meeting; Project 2 presentations and due date \\
  \textbf{December 9}  & Reading day; no class \\
\end{tabularx}

%==============================================================================
\section{Required Laptop and Software}

\textbf{Every student must bring a laptop to every class meeting.} Both meetings each week
involve live, hands-on work: reading and tracing code, running and repairing scripts,
inspecting variables and output, and building the week's studio deliverable. A student without
a working laptop cannot participate in the in-class activity that constitutes most of the
graded work in this course.

\begin{itemize}
  \item \textbf{Spyder IDE and Python 3.} Follow the current \emph{Spyder Installation and
        Course Environment Guide} posted on Blackboard.
        Week 1 verifies the Editor, Console, Variable Explorer, working directory, and course
        environment. Use the posted guide's version-specific instructions rather than an
        installation command found elsewhere.
  \item \textbf{Course libraries.} The course uses NumPy, pandas, Matplotlib, GeoPandas,
        SciPy, and scikit-learn. Libraries are introduced and verified when they are first
        needed; exact installation and troubleshooting instructions are provided in the setup
        guide and weekly materials.
  \item \textbf{Generative AI coding assistant.} The course will identify at least one
        no-cost option. Students are not required to purchase a paid AI service. Use of any
        assistant is governed by the Artificial Intelligence Policy below.
\end{itemize}

If your laptop cannot run Spyder, or if you do not have regular access to a laptop, contact
the instructor in the first week so an arrangement can be made before the graded studio work
begins.

%==============================================================================
\newpage
\section{Course Expectations}

The following guidelines are designed to facilitate a comfortable and productive learning
environment.

\begin{center}
\renewcommand{\arraystretch}{1.18}
\begin{tabularx}{\textwidth}{@{}>{\RaggedRight}X @{\hspace{1.2em}} >{\RaggedRight}X@{}}
\toprule
\textbf{Student Expectations} & \textbf{Instructor Expectations} \\
\midrule
Be punctual for all classes
  & Listen to and respect students' views \\
Bring a laptop to every class meeting, charged and able to run Spyder
  & Be available for questions at least 5 minutes before and after class \\
Read and review class notes prior to the lesson
  & Do not go beyond the allotted class time \\
Be active and engaged in class discussions and studio work
  & Accommodate differences in students' learning methodologies \\
Use laptops and devices for course work during class, not for unrelated activity
  & Create an active learning environment to facilitate student learning \\
Listen to and respect others
  & Reply to e-mails within 36 hours on weekdays and 48 hours on weekends \\
Complete all weekly studio assignments and both projects
  & Assign work that reflects the material covered in class \\
Submit all work before the stated deadline, with a complete AI-use record
  & Return feedback in time to inform the next week's work \\
\bottomrule
\end{tabularx}
\end{center}

\noindent\textit{If unable to meet the above expectations, contact the instructor ahead of
time.}

Note the deliberate difference from a lecture-only course: laptops are \emph{required} here,
not discouraged. The corresponding expectation is that the laptop is used for the course's own
work during class time.

%==============================================================================
\section{Attendance}

Each student has an educational obligation to attend all classes. Attendance will be checked
for every meeting. Because Wednesday studios build the graded deliverable in class with
instructor support present, an absence from a Meeting B is materially harder to recover from
than a missed lecture.

Arrival more than 10 minutes after the start of class is recorded as late; two late arrivals
count as one absence. Arrival more than 15 minutes late or departure more than 15 minutes
before the end of class counts as an absence. The first two absences carry no numeric
attendance penalty. This two-absence buffer is separate from the determination that an
absence is excused. Students who miss class remain responsible for notes, announcements, and
assigned work. Four or more unexcused absences will trigger a required conference about the
student's ability to complete the course successfully and, when timely, the option to
withdraw before the University deadline.

Per University policy, Barnes Center at the Arch staff do not provide medical excuse notes for
students. If Barnes Center staff determine it is medically necessary to remove a student from
classes, they will coordinate with the case management staff directly to provide absence
notification to faculty through Orange SUccess. For absences lasting less than 48 hours,
students are encouraged to discuss academic arrangements directly with their faculty.
Additional information may be found at Student Outreach and Retention:
\href{https://experience.syracuse.edu/soar/student-support/absence-notifications/}{Absence
Notifications}.

\subsection{University Expectation Regarding Attendance}

Attendance in classes is expected in all courses at Syracuse University. Students are expected
to arrive on campus in time to attend the first meeting of all classes for which they are
registered. Students who do not attend classes starting with the first scheduled meeting may
be academically withdrawn as not making progress toward degree by failure to attend.
Instructors set course-specific policies for absences from scheduled class meetings in their
syllabi. Students should also review the university's religious observance policy and make the
required arrangements at the beginning of each semester.

\subsection{Determination of Excused / Non-excused Absence}

Whether an absence is excused is determined by the instructor after reviewing the available
University notification or supporting information. Students should communicate promptly when
an emergency, religious observance, University-authorized activity, disability-related need,
or other serious circumstance affects attendance. An Orange SUccess absence notification
communicates that an absence occurred; it does not by itself determine whether the absence is
excused or waive missed course requirements. Approved accommodations and the University's
Religious Observances Policy remain in effect.

\subsection{Penalty for Unexcused Absences}

After the two-absence buffer, each additional unexcused absence results in a 1.5-point
deduction from the final course score. For example, three unexcused absences produce one
1.5-point deduction; the first two do not produce a numeric penalty.

%==============================================================================
\section{Weekly Studio Assignments}

A weekly assignment in this course \emph{is} that week's Wednesday studio handout, completed
and submitted as the graded artifact---not a second, separately designed deliverable layered
on top. Work begins in the Wednesday studio with instructor support in the room; the polished,
independently completed version is submitted afterward.

Every studio handout uses the same four-heading \texttt{\# \%\%} cell structure, and that
structure is the rubric skeleton:

\begin{center}
\texttt{QUESTION} \, / \, \texttt{INPUTS/ASSUMPTIONS} \, / \, \texttt{METHOD} \, / \,
\texttt{CHECKS-INTERPRET}
\end{center}

\subsection{Submission Rhythm and Deadlines}

\begin{itemize}
  \item \textbf{Standard rule:} due before the \emph{following} Monday's Meeting A, at
        \textbf{9:00 am}. That Monday then opens with a 3--5 minute debrief of the most common
        error before new content begins. Weeks 1, 3, 4, 5, 6, 9, 10, 11, and 12 follow this
        rule unchanged.
  \item \textbf{Three exceptions}, where the following Monday has no meeting to submit before
        and no session in which to debrief. Each is due the \textbf{Saturday before, 11:59
        pm}, and each debrief moves to the next session that actually meets:
        \begin{itemize}
          \item \textbf{Week 2} (studio Wed Sep 2) $\rightarrow$ due Saturday, September 5.
                The following Monday is Labor Day.
          \item \textbf{Week 7} (studio Wed Oct 7) $\rightarrow$ due Saturday, October 10.
                The following Monday is Fall Break.
          \item \textbf{Week 14} (studio Wed Dec 2) $\rightarrow$ due Saturday, December 5.
                The following Monday is Project 2's own due date and presentation day, so this
                gap preserves real runway between the last checkpoint and the final
                submission.
        \end{itemize}
  \item \textbf{One shifted deadline:} Week 13's assignment (studio Wed Nov 18) is due Monday,
        November 30 at \textbf{12:45 pm}, the start of Meeting A, because Thanksgiving break
        intervenes.
  \item \textbf{Weeks 8 and 15} carry no separate weekly assignment---their studio time
        \emph{is} the project deliverable. This is why the weekly bucket is built from 13
        graded weeks (1--7 and 9--14), not 15.
\end{itemize}

\subsection{Late Policy}

15\% off within 24 hours late; zero credit after 48 hours, barring a prior arrangement made in
advance of the deadline. Technical difficulties such as a poor internet connection are not
considered an emergency justifying an extension. Students are required to keep an electronic
record of each submission and to confirm that Blackboard returned a submission receipt.
Incomplete submission will be considered non-submission.

\subsection{Lowest Two Dropped}

The two lowest weekly studio grades are dropped, so one bad week is not catastrophic.

\subsection{The Transfer Safeguard}

Starting in Week 4, every submission's \texttt{CHECKS-INTERPRET} section must include at least
one self-generated check applied to a slightly different input than the one worked through
live in class---a different row, case, parameter value, or small scenario tweak---with 2--4
\texttt{assert} statements the student writes and runs before submitting. Copying the in-class
values alone does not satisfy this requirement.

\subsection{Individual Work}

All weekly assignments are individual. Discussing main concepts or general approaches is
acceptable; joint problem-solving, sharing files, or submitting work substantially derived
from another student's is not. Both students implicated---the one who provided the work and
the one who replicated it---share equal responsibility for a breach of academic integrity
under University policy. Using or disseminating solutions from prior years' assignments is
likewise a violation.

%==============================================================================
\section{Projects}

This course has \textbf{no midterm or final exam}. In their place are two cumulative projects,
each with milestones that \emph{are} the ordinary weekly studio task for that week rather than
additional artifacts layered on top. From Week 5 onward, students apply each week's studio
skill directly to their own project data.

\subsection{Project 1 --- Evidence Before Design (assigned September 14; due October 14)}

A building-performance and site data story built from the student's own chosen dataset.
Deliverables: one or two \texttt{.py} files (\texttt{analysis.py} and, if needed,
\texttt{evidence.py}); one cleaned CSV and one grouped-summary CSV; two static PNG figures
plus either one map PNG or one animation GIF; a 600--800-word illustrated evidence memo; a
README, data dictionary, and AI-use record; and a six-slide pin-up, using the provided
template, presented in the Week 8 clinic.

The six-slide pin-up takes place during class on Wednesday, October 14. The complete digital
submission is due in Blackboard by \textbf{11:59 pm Eastern Time on Wednesday, October 14,
2026}. Week 7's map remains a graded weekly studio. For the final Project 1 evidence
extension, students may submit either that map or the Week 6 animation; students choosing the
animation still complete the Week 7 spatial studio as weekly course work.

\begin{center}
\begin{tabular}{@{}lr@{}}
\toprule
\textbf{Project 1 rubric} & \textbf{Weight} \\
\midrule
Question, provenance, cleaning, checks, and cleaned export & 20\% \\
Grouped evidence and defensible outlier decision           & 20\% \\
Two purposeful static figures                              & 20\% \\
Chosen extension: map OR animation                         & 15\% \\
Interpretation, limitations, and bounded recommendation     & 15\% \\
Reproducibility, code clarity, README, and AI-use record    & 10\% \\
\midrule
\textbf{Total} & \textbf{100\%} \\
\bottomrule
\end{tabular}
\end{center}

\textbf{Milestones} (each is that week's regular studio assignment): Week 5 --- dataset,
research question, and provenance, via the grouping and outlier analysis; Week 6 ---
visualization checkpoint; Week 7 --- spatial studio plus the compiled draft of Parts I--III
already produced in Weeks 4--6.

\subsection{Project 2 --- Predict, Optimize, Decide (assigned October 19; due December 7)}

Performance-informed design alternatives, integrating a fitted prediction model into an
optimization on the student's own problem. Deliverables: a one-page formulation and sensitivity
memo; two or three \texttt{.py} files (\texttt{model\_and\_evaluate.py},
\texttt{optimize\_and\_integrate.py}, and an optional helper); one model metric table, one
diagnostic figure, and one model card; one primary optimization result, one compact solver
comparison, one hand-verified constraint table, and one solver card; four purposeful figures
(sensitivity/design space, model diagnostic, solver comparison, and baseline versus
recommended design); a 1,000--1,200-word report including the integration step and its
limitations; a README and AI-use record; and a five-minute presentation: a four-minute
decision brief plus a one-minute evidence/code trace, hard-capped at five minutes per
student in the Week 15 session.

The presentations take place during class on Monday, December 7. The final 15 minutes are
reserved for selective source-deck synthesis. If enrollment exceeds 12, presentations use
parallel gallery/code-review rounds before the class reconvenes for that synthesis.
The complete digital submission is due in Blackboard by \textbf{11:59 pm Eastern Time on
Monday, December 7, 2026}.

\begin{center}
\begin{tabular}{@{}lr@{}}
\toprule
\textbf{Project 2 rubric} & \textbf{Weight} \\
\midrule
Formulation, family justification, and sensitivity evidence & 15\% \\
Evaluated prediction, diagnostic, and model card            & 20\% \\
Primary optimization and constraint verification            & 20\% \\
Compact solver / baseline comparison                        & 15\% \\
Predict--optimize--decide integration                       & 20\% \\
Communication, reproducibility, limitations, and AI-use record & 10\% \\
\midrule
\textbf{Total} & \textbf{100\%} \\
\bottomrule
\end{tabular}
\end{center}

\textbf{Milestones} (each is that week's regular studio assignment): Week 10 --- topic,
dataset, and pairing proposal via the \texttt{linprog} formulation of the student's own linear
sub-problem; Week 11 --- nonlinear formulation; Week 12 --- heuristic-search track; Week 13 ---
baseline model, fitted, evaluated, and leakage-audited; Week 14 --- evaluated prediction and
model card.

%==============================================================================
\section{Grades}

The course evaluation is built from the weekly studio assignments and the two projects. There
is no separate quiz or exam category, and no separately weighted participation line.

\begin{center}
\begin{tabular}{@{}lr@{}}
\toprule
\textbf{Component} & \textbf{Weight} \\
\midrule
Weekly studio assignments (13 graded weeks; lowest 2 dropped) & 40\% \\
Project 1 --- Evidence Before Design (due Week 8)              & 25\% \\
Project 2 --- Predict, Optimize, Decide (due Week 15)          & 35\% \\
\midrule
\textbf{Total} & \textbf{100\%} \\
\bottomrule
\end{tabular}
\end{center}

Two policies are deliberately \emph{not} weighted line items, and instead act on the composite
score:

\begin{itemize}
  \item \textbf{Attendance and professionalism} applies as a flat deduction on the final
        composite score, as described under \textit{Penalty for Unexcused Absences}.
  \item \textbf{AI-use disclosure} is graded inside each weekly assignment's and each
        project's own rubric, and additionally functions as a compliance gate: a missing,
        incomplete, or fabricated disclosure caps that submission's score at \textbf{70\%}
        regardless of technical quality.
\end{itemize}

A zero-prior-experience cohort meeting twice weekly for 80 minutes benefits most from a short,
frequent, recoverable feedback loop rather than a few high-stakes tests. Hence 40\% is earned
in small weekly steps with the two lowest dropped, and the remaining 60\% rewards cumulative
integration. If students are expected to use AI coding assistants throughout, a
memorized-syntax exam measures the wrong thing; the weekly assignments provide the frequent,
individually attributable check instead.

\subsection{Grade Scale}

A grade ``A'' merits work that pushes to levels of excellence beyond the expectations of the
class. A grade ``B'' merits very good work that fulfills the expectations of the class. A
grade ``C'' merits work that fulfills the minimum requirement for the class. Work that does
not fulfill the minimum expectations of this graduate-level course will receive a grade
``F.'' Syracuse University does not permit D or D$-$ grades in graduate courses; ARC 500 uses
the graduate grading standard for all enrolled students.

Numerical scale (all numbers round to the nearest whole; e.g.\ 68.52 = 69 = C while 68.33 = 68
= C$-$):

\begin{center}
\begin{tabular}{@{}llcll@{}}
\toprule
\textbf{Percentage} & \textbf{Letter} & \hspace{1.5em} & \textbf{Percentage} & \textbf{Letter} \\
\midrule
93 and above & A    & & 77--80 & B$-$ \\
89--92 & A$-$ & & 73--76 & C+  \\
85--88 & B+   & & 69--72 & C   \\
81--84 & B    & & 65--68 & C$-$ \\
Below 65 & F & & & \\
\bottomrule
\end{tabular}
\end{center}

%==============================================================================
\section{Course Materials}

\textbf{There is no required textbook to purchase for this course.} Class notes, the weekly
studio handouts, and the posted setup guide are the primary materials, and every graded
artifact is fully specified in those materials on its own.

A \textbf{companion text} written for this course is available as a free, living document:
\emph{Introduction to Programming with Python} (J. Chun), an architecture-framed Python text
that uses the same room, facade, beam, and grid examples as the lecture material. Nothing in
this course depends on a student reading it; it is offered as optional deeper reading where a
week's lecture cites a chapter. One deliberate divergence: the companion text installs and
edits in VS Code, while this course uses Spyder throughout---follow the course's Spyder
instructions where the two differ.

As of August 2026, the optional weekly references are: Week 1, Chapters 1 and 8; Week 2,
Chapter 2; Week 3, Chapters 3--6; Week 4, Chapters 10--12; Week 5, Chapter 12; Week 6,
Chapter 13; Week 7, Chapters 15--16; Week 9, Chapter 11; Week 11, Chapter 18; Week 12,
Chapter 20; Week 13, Chapter 19; and Week 14, Chapter 21. Weeks 8, 10, and 15 have no matching
assigned chapter. Chapter 22 (Neural Networks) is beyond the scope of ARC 500. Because this is
a living document, the week-specific link posted on Blackboard is the reference of record if
chapter numbering changes.

\subsection{Reference Documentation}

Official documentation is the reference of record for every library used, and learning to read
it is itself a course skill:

\begin{itemize}
  \item Python: \url{https://docs.python.org/3/}
  \item Spyder: \url{https://docs.spyder-ide.org/current/}
  \item NumPy: \url{https://numpy.org/doc/stable/}
  \item pandas: \url{https://pandas.pydata.org/docs/}
  \item Matplotlib: \url{https://matplotlib.org/stable/}
  \item GeoPandas: \url{https://geopandas.org/en/stable/}
  \item SciPy \texttt{optimize}: \url{https://docs.scipy.org/doc/scipy/reference/optimize.html}
  \item scikit-learn: \url{https://scikit-learn.org/stable/}
\end{itemize}

%==============================================================================
\section{Course Role in the Program}

\ProgramRole

%==============================================================================
\section{Course Contents}

\begin{enumerate}
  \item \textbf{Computational thinking and the Python execution model}
        \begin{itemize}
          \item Decomposition, abstraction, and automation
          \item Problem formulation before syntax
          \item Scripts, statements, comments; syntax, runtime, and logic error
          \item The generative-AI-as-collaborator protocol
          \item The Spyder IDE: Editor, Console, Variable Explorer, kernel, tracebacks
        \end{itemize}
  \item \textbf{Values, types, and precise expression}
        \begin{itemize}
          \item \texttt{int}, \texttt{float}, \texttt{str}, \texttt{bool}, \texttt{None}
          \item Naming with units; operator precedence; unit conversion
          \item f-strings and formatted output; floating-point display
          \item Comparisons; \texttt{assert} as a first correctness check
        \end{itemize}
  \item \textbf{Decisions, repetition, and reusable logic}
        \begin{itemize}
          \item \texttt{if} / \texttt{elif} / \texttt{else}; boundary cases; compound conditions
          \item \texttt{for} loops over a list; the accumulator pattern
          \item Lists and a minimal dictionary lookup
          \item \texttt{def}, parameters, defaults, \texttt{return}; guard clauses
        \end{itemize}
  \item \textbf{The toolchain and data as evidence}
        \begin{itemize}
          \item NumPy arrays: shape, dtype, indexing, vectorized arithmetic
          \item pandas Series and DataFrame; the trustworthiness interview
          \item Data dictionaries, provenance, and NaN as a modeling decision
          \item Docstrings and type hints; filtering, derived columns, and CSV export
        \end{itemize}
  \item \textbf{Grouping, outliers, and defensible summaries}
        \begin{itemize}
          \item Split--apply--combine; \texttt{groupby()}, multiple aggregations, sorting
          \item IQR versus $z$-score outlier rules and where they disagree
          \item Measurement error versus genuine extreme versus a missing variable
          \item One key-based merge
        \end{itemize}
  \item \textbf{Visual reasoning: static and dynamic evidence}
        \begin{itemize}
          \item Matching a question to an encoding: distribution, comparison, relationship,
                change over a sequence
          \item The \texttt{fig, ax} pattern; labeled axes with units; legends
          \item Annotating the one point that carries the argument
          \item Dynamic evidence: \texttt{FuncAnimation} saved as a shareable GIF
        \end{itemize}
  \item \textbf{Spatial data, scale, and maps}
        \begin{itemize}
          \item Geometry versus attributes; point, line, polygon; GeoDataFrame
          \item CRS and EPSG codes; why area and length require a projected CRS
          \item Buffer and spatial join; a properly labeled static choropleth
        \end{itemize}
  \item \textbf{Parameter sweeps and sensitivity}
        \begin{itemize}
          \item Vectorized operations versus explicit loops
          \item \texttt{arange} / \texttt{linspace}; boolean masks; \texttt{meshgrid}
          \item Reading a swept surface: which input moves the outcome most
        \end{itemize}
  \item \textbf{Optimization I: formulation and exact linear methods}
        \begin{itemize}
          \item Design variables, objective, constraints, bounds, feasible region
          \item Standard form; why a linear problem is solved exactly rather than searched
          \item \texttt{scipy.optimize.linprog}; hand-verified constraint tables
        \end{itemize}
  \item \textbf{Optimization II: nonlinear and gradient-based search}
        \begin{itemize}
          \item \texttt{scipy.optimize.minimize}; initial guess and bounds
          \item Local versus global optima; multi-start comparison
        \end{itemize}
  \item \textbf{Optimization III: heuristic and metaheuristic search}
        \begin{itemize}
          \item Discrete and combinatorial choices; non-smooth objectives
          \item Random search, hill-climbing with restarts, genetic algorithms
                (selection, crossover, mutation), simulated annealing
          \item Comparison against a brute-force baseline, in result and runtime
        \end{itemize}
  \item \textbf{Machine learning I: evaluated prediction (regression)}
        \begin{itemize}
          \item Features versus target; supervised regression; \texttt{train\_test\_split}
          \item Why a naive baseline comes first; MAE, RMSE, $R^2$; residual plots
          \item Overfitting, underfitting, and data leakage---found and fixed in AI-drafted code
        \end{itemize}
  \item \textbf{Machine learning II: classification, model cards, and decisions}
        \begin{itemize}
          \item Classification as a decision rule; the confusion matrix
          \item Precision versus recall; choosing a threshold from real consequences
          \item \texttt{cross\_val\_score}; the model card as a deliverable
        \end{itemize}
  \item \textbf{Synthesis: predict, optimize, decide}
        \begin{itemize}
          \item A fitted model as a surrogate objective inside an optimization
          \item Reproducibility and the AI-use record
          \item Stating a bounded recommendation with its limitations
        \end{itemize}
\end{enumerate}

%==============================================================================
\section{Blackboard}

This class will use Blackboard for all materials relevant to the course, including handouts,
studio files, assignment specifications, solutions, links to external materials, and grades.
Blackboard should be visited frequently to obtain the latest versions of class notes,
assignments, and reference materials. All deadline times for submitting materials on
Blackboard are Eastern Time.

For Blackboard or account support, contact Information Technology Services at
\href{mailto:help@syr.edu}{help@syr.edu}, 315.443.2677, or the ITS Service Center in 1-227
CST, Life Sciences Complex.

\subsection{Blackboard Discussions.}

Use Blackboard Discussions for questions likely to help the class. Do not post a complete
working solution to a current assignment; posting the exact error message and the shortest
relevant code excerpt, together with what you already tried, is encouraged.

%==============================================================================
\section{Academic Drop Deadline}

The academic/financial deadline to drop a class is \DropDeadline{}. This is also the
grading-option deadline and the religious-observance notification deadline. The University
withdrawal deadline is \WithdrawalDeadline{}. A course withdrawal after the drop deadline and
by the withdrawal deadline remains on the transcript with the grading symbol \texttt{WD}.

%==============================================================================
\section{Syracuse University Policies}

The following required statements reflect Syracuse University's current syllabus guidance.

\subsection{Disability Statement}

Syracuse University values diversity and inclusion and is committed to a climate of mutual
respect and full participation. There may be aspects of the instruction or design of this
course that result in barriers to your inclusion and full participation. I invite students to
meet with me to discuss strategies and/or accommodations (academic adjustments) that may be
essential to your success in collaboration with the
\href{https://disabilityresources.syr.edu/}{Center for Disability Resources} (CDR).

To discuss disability-accommodations or register with CDR, visit
\href{https://disabilityresources.syr.edu/}{disabilityresources.syr.edu}, or contact the
center at 315.443.4498 or \href{mailto:disabilityresources@syr.edu}{disabilityresources@syr.edu}
for more detailed information. CDR is responsible for coordinating disability-related
academic accommodations and will work with you to develop an access plan. Since accommodations
may require early planning and generally are not provided retroactively, please contact CDR as
soon as possible.

If any aspect of the required in-class laptop work presents an access barrier, please contact
me in the first week so an arrangement can be made before the graded studio work begins.

\subsection{Academic Integrity Policy}

As a preeminent and inclusive student-focused research institution, Syracuse University
places academic integrity at the forefront of learning and considers it a core value and
guiding pillar of education. The University's
\href{https://policies.syr.edu/policies/academic-rules-student-responsibilities-and-services/academic-integrity-policy/}{Academic Integrity Policy}
provides guidelines for completing academic work with integrity.

Students are required to meet both course-specific and Universitywide academic integrity
expectations, such as crediting your sources, doing your own work, communicating honestly and
supporting academic integrity.

Upholding academic integrity includes the protection of faculty intellectual property.
Students should not upload, distribute or share an instructor's course materials, including
presentations, assignments, exams or other evaluative materials, without permission.

Using websites that charge fees or require uploading of course material (e.g., Chegg, Course
Hero) to obtain exam solutions or assignments completed by others and presenting them as your
own violates academic integrity expectations in this course and may be classified as a Level
3 violation.

All academic integrity expectations that apply to in-person assignments, quizzes and exams
also apply online. Students found in violation of the policy are subject to grade sanctions
determined by the course instructor and non-grade sanctions determined by the school or
college offering the course. Students may not drop or withdraw from courses in which they face
a suspected violation. Any established violation in this course may result in course failure
regardless of violation level.

\subsection{Artificial Intelligence Policy}

\textbf{Full Use with Disclosure and Citation.} Based on the assignments and specified
learning outcomes, the full use of artificial intelligence as a tool, with disclosure and
citation, is permitted in this course.

\textbf{Generative AI is permitted in this course, and using it well is one of the course's
learning goals.} This is a deliberate, course-specific policy: ARC 500 treats an AI coding
assistant as a collaborator whose output must be verified, and much of the course's value lies
in learning to formulate problems precisely, evaluate what an assistant returns, and repair
what it gets wrong. Several assignments explicitly require auditing AI-drafted code and fixing
a real defect in it.

Permission is conditional on all of the following. Each is a requirement, not a suggestion:

\begin{enumerate}
  \item \textbf{You must be able to explain every line you submit.} You are responsible for
        being able to trace, test, modify, and explain any submitted line on request. ``The
        assistant wrote it'' is not an explanation, and inability to explain your own
        submission is treated as an academic integrity concern regardless of whether the code
        runs.
  \item \textbf{You must keep and submit an AI-use record.} Every weekly assignment and both
        projects require a disclosure stating which tool was used, what was asked of it, what
        it returned, what you kept, what you changed or rejected and why, and how you tested
        each retained suggestion. A missing, incomplete, or fabricated disclosure
        \textbf{caps that submission at 70\%} regardless of technical quality. A fabricated
        disclosure---claiming you did not use a tool that you did use, or inventing a testing
        process you did not perform---is a dishonesty violation, not merely a
        documentation lapse.
  \item \textbf{You must verify, not trust.} AI-generated code in this course is treated as a
        draft to be checked. Verification means running it, testing it against known values
        with \texttt{assert}, and confirming that the numbers, units, and interpretation are
        correct. Submitting unverified output that happens to run is not acceptable work.
  \item \textbf{AI does not substitute for the reasoning being assessed.} Where an assignment
        asks you to interpret a result, defend a decision, state a limitation, or write a
        memo, that reasoning must be your own. Using a tool to draft prose you then adopt
        without independent judgment defeats the purpose of that component and will be graded
        accordingly.
  \item \textbf{AI use does not license collaboration.} Assignments remain individual work.
        Sharing prompts, transcripts, or generated solutions with other students is the same
        violation as sharing hand-written code.
  \item \textbf{Do not submit confidential or third-party material to an AI tool.} Do not
        upload course materials that are the instructor's intellectual property, other
        students' work, or any data you are not permitted to share.
\end{enumerate}

If you are unsure whether a particular tool or feature is considered generative AI, or whether
a particular use is acceptable, ask the instructor \emph{before} using it. Asking is always
acceptable and is never penalized.

%==============================================================================
\section{Faith Tradition Observances}

Syracuse University's Religious Observances Policy recognizes the diversity of faiths
represented in the campus community and protects the rights of students, faculty and staff to
observe religious holy days according to their traditions. Under the policy, students are given
an opportunity to make up any examination, study, or work requirements that may be missed due
to a religious observance, provided they notify their instructors no later than the academic
drop deadline.

For any observances occurring before the academic drop deadline, students must notify faculty
at least two academic days in advance. Students register their observances using MySlice.

%==============================================================================
\section{Discrimination or Harassment}

Syracuse University does not discriminate and prohibits harassment or discrimination related
to any protected category including creed, ethnicity, citizenship status, reproductive health
decisions, national origin, sex, gender, pregnancy, disability, marital status, age, race,
color, veteran status, military status, religion, sexual orientation, domestic violence
status, genetic information, gender identity and/or gender expression, shared common ancestry
or any other status protected by applicable law.

Any complaint of discrimination or harassment related to any of these protected bases should
be reported to Sheila Johnson-Willis, associate vice president and chief equal opportunity and
Title IX officer in the Office of Equal Opportunity, Inclusion and Resolution Services (621
Skytop Road, Suite 1001, Syracuse, NY 13244;
\href{mailto:equalopp@syr.edu}{equalopp@syr.edu}; 315.443.4018). She is responsible for
coordinating compliance efforts under various laws including Titles VI, VII, IX and Section
504 of the Rehabilitation Act.

%==============================================================================
\section{Use of Class Materials}

Original class materials (lectures, notes, syllabus, study guides, handouts, assignments,
studio files, projects, programs, and tools) of this course are the intellectual property of
the course instructor (Junho Chun) and are protected by U.S. copyright law and by University
policy. The course instructor is the exclusive owner of the copyright in those materials.
Students may take notes and make copies of course materials for their own use. However,
students should not provide these materials to other parties (i.e., websites, social media, or
any student not registered or enrolled in this course) without written permission from the
instructor. Doing so is a violation of intellectual property law and of the
\href{https://policies.syr.edu/policies/academic-rules-student-responsibilities-and-services/code-of-student-conduct/}{Code
of Student Conduct}.

This includes uploading course materials into a generative AI tool that retains or trains on
submitted content.

%==============================================================================
\section{Student Work}

I intend to use academic work that you complete this semester in subsequent semesters for
educational purposes. Before using your work for that purpose, I will either get your written
permission or render the work anonymous by removing identifying material. The School of
Architecture may also request digital versions of selected course work for its Student Works
Archive. Selection criteria, file-naming protocols, and required formats will be explained
during the semester; any required archive package will be evaluated within the applicable
project's professional-submission and reproducibility criteria.

%==============================================================================
\section{University Email Policy}

Syracuse University email is the official channel for University communication. Students are
responsible for monitoring their syr.edu accounts. Faculty and staff must use their official
University email addresses when communicating with students about course matters; sending
student records to a non-syr.edu address may create a FERPA violation. See the
\href{https://policies.syr.edu/policies/information-technology/e-mail-policy/}{Syracuse
University Email Policy}.

%==============================================================================
\section{Use of Turnitin}

This class will use the plagiarism detection and prevention system Turnitin for its written
deliverables (the Project 1 evidence memo and the Project 2 report). Students will have the
option to submit their writing to Turnitin to check that all sources have been properly
acknowledged and cited before final submission. The instructor will also submit these
documents to Turnitin, which compares submitted documents against documents on the Internet
and against student papers submitted to Turnitin at Syracuse University and at other colleges
and universities. Students' knowledge of the subject matter of this course and their writing
level and style will be considered in interpreting the originality of a report. Keep in mind
that all written work you submit for this class will become part of the Turnitin.com reference
database solely for the purpose of detecting plagiarism of such papers.

%==============================================================================
\newpage
\section{Course Schedule}

\newcolumntype{K}{>{\RaggedRight\footnotesize}p{0.055\textwidth}}
\newcolumntype{J}{>{\RaggedRight\footnotesize}p{0.275\textwidth}}
\newcolumntype{Q}{>{\RaggedRight\footnotesize}p{0.30\textwidth}}

\renewcommand{\arraystretch}{1.15}
\begin{longtable}{@{}K J J Q@{}}
\toprule
\textbf{Wk} & \textbf{Monday (Meeting A)} & \textbf{Wednesday (Meeting B)} &
\textbf{Assignment \& Note} \\
\midrule
\endfirsthead
\multicolumn{4}{@{}l}{\small\itshape Course schedule, continued} \\
\toprule
\textbf{Wk} & \textbf{Monday (Meeting A)} & \textbf{Wednesday (Meeting B)} &
\textbf{Assignment \& Note} \\
\midrule
\endhead
\midrule
\multicolumn{4}{@{}r@{}}{\small\itshape continued on next page} \\
\endfoot
\bottomrule
\endlastfoot

1 & \textbf{Aug 24} \newline Computational thinking and the Python execution model
  & \textbf{Aug 26} \newline Spyder installation and first program
  & Week 1 studio: trace and repair a starter script. Due Mon Aug 31, 9:00 am. \\

2 & \textbf{Aug 31} \newline Values, types, and precise expression
  & \textbf{Sep 2} \newline Architectural quantity studio
  & Week 2 studio: unit-aware calculator with \texttt{assert} checks.
    \textbf{Due Sat Sep 5, 11:59 pm} (Labor Day exception). \\

3 & \textbf{Sep 7} \newline \textit{No class: Labor Day.} Asynchronous primer
    ($\sim$40 min) on decisions and functions, completed before Wednesday
  & \textbf{Sep 9} \newline Design-rule and function studio
  & Week 3 studio: validated function, loop, dictionary lookup, \texttt{assert} tests.
    Due Mon Sep 14, 9:00 am. \\

4 & \textbf{Sep 14} \newline Meet the toolchain, then data as evidence
  & \textbf{Sep 16} \newline Trustworthiness studio
  & \textbf{Project 1 assigned.} Week 4 studio: cleaned dataset, data dictionary,
    drop-justification memo, plus one NumPy-derived column plotted as a single
    Matplotlib scatter. Due Mon Sep 21, 9:00 am. \\

5 & \textbf{Sep 21} \newline Grouping, outliers, and defensible summaries
  & \textbf{Sep 23} \newline Grouping and outlier studio
  & Week 5 studio (= \textbf{P1 milestone}: dataset, question, provenance):
    grouped table, dual-rule outlier flags, decision sentences.
    Due Mon Sep 28, 9:00 am. \\

6 & \textbf{Sep 28} \newline Visual reasoning: static and dynamic evidence
  & \textbf{Sep 30} \newline Visual evidence studio
  & Week 6 studio (= \textbf{P1 visualization milestone}): three static figures and one
    captioned animated GIF. Due Mon Oct 5, 9:00 am. \\

7 & \textbf{Oct 5} \newline Spatial data, scale, and maps
  & \textbf{Oct 7} \newline Spatial join and choropleth studio
  & Week 7 studio (= \textbf{P1 spatial layer + compiled draft}): GeoPandas script,
    choropleth, CRS justification. \textbf{Due Sat Oct 10, 11:59 pm} (Fall Break
    exception). \\

8 & \textbf{Oct 12} \newline \textit{No class: Fall Break.}
  & \textbf{Oct 14} \newline Project 1 studio and evidence clinic
  & \textbf{Project 1 pin-up in class; digital package due 11:59 pm.} No separate weekly
    assignment; the clinic, pin-up, and submission are this week's work. \\

9 & \textbf{Oct 19} \newline Parameter sweeps and sensitivity
  & \textbf{Oct 21} \newline Sweep and heatmap studio
  & \textbf{Project 2 assigned.} Week 9 studio: two-parameter sweep, heatmap, sensitivity
    finding, vectorized-vs-loop \texttt{assert}. Due Mon Oct 26, 9:00 am. \\

10 & \textbf{Oct 26} \newline Optimization I: formulation and exact linear methods
  & \textbf{Oct 28} \newline \texttt{linprog} studio
  & Week 10 studio and \textbf{P2 proposal}: formulate your own linear sub-problem with
    \texttt{linprog} and submit a hand-verified constraint table.
    Due Mon Nov 2, 9:00 am. \\

11 & \textbf{Nov 2} \newline Optimization II: nonlinear and gradient-based search
  & \textbf{Nov 4} \newline Multi-start optimization studio
  & Week 11 studio (= \textbf{P2 nonlinear checkpoint}): \texttt{minimize} run plus a
    multi-start comparison write-up. Due Mon Nov 9, 9:00 am. \\

12 & \textbf{Nov 9} \newline Optimization III: heuristic and metaheuristic search
  & \textbf{Nov 11} \newline Genetic algorithm studio
  & Week 12 studio (= \textbf{P2 heuristic checkpoint}): GA or SA implementation versus a
    brute-force baseline. Due Mon Nov 16, 9:00 am. \\

13 & \textbf{Nov 16} \newline Machine learning I: evaluated prediction (regression)
  & \textbf{Nov 18} \newline Regression studio
  & Week 13 studio (= \textbf{P2 baseline-model checkpoint}): baseline-vs-model metric
    table, residual plot, and a fixed leakage bug in AI-drafted code.
    \textbf{Due Mon Nov 30, 12:45 pm}, at the start of Meeting A (Thanksgiving shift). \\

-- & \textbf{Nov 23} \newline \textit{No class: Thanksgiving Break}
  & \textbf{Nov 25} \newline \textit{No class: Thanksgiving Break}
  & Unnumbered break week; no new content. \\

14 & \textbf{Nov 30} \newline Machine learning II: classification, model cards, and
     decisions
  & \textbf{Dec 2} \newline Classification / regression studio (track-dependent)
  & Week 14 studio (= \textbf{P2 model-card checkpoint}): track-appropriate model update,
    evaluation artifact, one-page model card. \textbf{Due Sat Dec 5, 11:59 pm.} Last week
    new skills are introduced. \\

15 & \textbf{Dec 7} \newline Capstone clinic, synthesis, and in-class presentations
     \textit{(the only meeting this week)}
  & \textbf{Dec 9} \newline \textit{No class: reading day}
  & \textbf{Project 2 presentations in class; digital package due 11:59 pm.} No separate
    weekly assignment. \\

\end{longtable}

\vspace{2pt}
{\footnotesize\itshape * Schedule subject to change. Any change will be announced in class and
posted on Blackboard.}

\vspace{10pt}
{\color{rulegray}\hrule height 0.4pt}
\vspace{4pt}
{\footnotesize\RaggedRight Calendar notes: classes begin Monday, August 24, 2026. Labor Day (Mon Sep 7)
removes Week 3's Meeting A. Fall Break (Mon Oct 12 and Tue Oct 13) removes Week 8's Meeting A;
Tuesday is not a class day for this Monday/Wednesday course. Thanksgiving Break runs Sun Nov
22 through Sun Nov 29 and removes both meetings that week, which is treated as an unnumbered
break between Weeks 13 and 14 so the numbered sequence stays 1--15. The last ARC 500 meeting
is Monday, December 7; December 9 is a reading day.\par}

\end{document}
